\section{The normalizer of the label subgroup}
\label{nor:section}

By Proposition~\ref{hyp:finite-out}, $\Out(\Pi)$ is finite.
It contains the two realized subgroups $\mathsf A$ and
$\mathsf B_\ell$, which are not conjugate by
Proposition~\ref{ext:nonconjugate}.  To apply the Sylow criterion,
we determine the normalizer of $\mathsf A$.  Conjugation by a
representative of a normalizer element preserves the full preimage
$\mathcal E_2$ of $\mathsf A$.  We will show that it preserves $W$
and that its restriction to $W$ is an inner automorphism followed by
an automorphism induced by a simplicial symmetry of the infinite nerve
$\widehat{\mathcal L}$; the label relations will then force that
symmetry to be a deck transformation.  The rigidity proof uses finite
supports and the topology of links.

\subsection{Rigidity of the Coxeter generators}

For a simplicial complex $X$, write $W(X)$ for the right-angled
Coxeter group on its graph.  In the subscript of a special
subgroup, a subcomplex denotes its vertex set.  For a graph
$\mathcal G$, the notation $\st_{\mathcal G}(u)$ means the
closed neighbor set consisting of $u$ and all its neighbors,
and $\lk_{\mathcal G}(u)$ means the neighbor set.  The following
facts are standard for finite graphs
\cite[Lemmas~2.4--2.7]{GutierrezPiggottRuane2012}.  The proof
records why they also hold for an arbitrary vertex set.

\begin{lemma}\label{nor:coxeter-facts}
Let $W(\mathcal G)$ be the right-angled Coxeter group on a simple
graph $\mathcal G$ with vertex set $V$.
\begin{enumerate}
\item For $u\in V$,
\[
 C_{W(\mathcal G)}(u)=W_{\st_{\mathcal G}(u)}.
\]
\item Every finite-order element is conjugate to a product of
distinct generators supported on a clique, with the empty
product allowed.
\item The center is the special subgroup on the universal
vertices:
\[
 Z(W(\mathcal G))
 =W_{\{v\in V:\ v\text{ is adjacent to every }u\ne v\}}.
\]
\end{enumerate}
\end{lemma}

\begin{proof}
Put $L=\lk_{\mathcal G}(u)$.  The presentations and the
special-subgroup injections of Lemma~\ref{ext:support} give
\[
 W(\mathcal G)
 =W_{V\setminus\{u\}}*_{W_L}
       \bigl(W_L\times\langle u\rangle\bigr).
\]
If $u$ is adjacent to every other vertex, the second factor is
the whole group and the centralizer assertion is immediate.
Otherwise consider the Bass--Serre tree of this amalgam.  The
element $u$ fixes the vertex of the second factor.  It fixes no
edge incident to that vertex: every such edge stabilizer is a
conjugate of $W_L$ within $W_L\times\langle u\rangle$, hence is
$W_L$, which does not contain $u$.  The fixed subtree of $u$
is therefore that single vertex.  Its centralizer must
stabilize this vertex and so lies in the second factor.  That
factor centralizes $u$, proving the first assertion.

For the second assertion, a finite-order element lies in a
special subgroup on a finite supporting set by
Lemma~\ref{ext:support}.  Work in that finite induced graph and
induct on its number of vertices.  A finite-order element in
the displayed amalgam fixes a vertex of its Bass--Serre tree,
so is conjugate into one of the factors.  The first factor has
fewer generators.  In the second factor its projection to
$W_L$ has finite order; induction conjugates that projection
to a product on a clique in $L$.  Multiplying by the possible
$u$-factor still gives a clique.  This also covers the
degenerate direct-product case, and starts with the trivial
group.

Finally, an element is central exactly when it belongs to all
the centralizers in the first assertion.  By
Lemma~\ref{ext:support}, its support must then lie in
$\bigcap_{u\in V}\st_{\mathcal G}(u)$, the set of universal
vertices.  Every generator in that set is central.  This gives
the asserted equality, also for an infinite graph.
\end{proof}

\begin{lemma}\label{nor:generator-lines}
Every automorphism of $W$ induces a permutation of the
coordinate vectors in
\[
 W_{\mathrm{ab}}
   =\bigoplus_{u\in\widehat{\mathcal V}}\Ftwo\,\mathbf e_u.
\]
The resulting vertex permutation is a simplicial automorphism
of $\widehat{\mathcal L}$.  After composing with the inverse of
that generator permutation, the original automorphism sends
each Coxeter generator to a conjugate of itself.
\end{lemma}

\begin{proof}
Let $\mathcal T\subseteq W_{\mathrm{ab}}$ consist of the
vectors represented by finite-order elements of $W$.
Lemma~\ref{nor:coxeter-facts} shows that these are exactly the
vectors supported on cliques, including the zero vector: the
converse holds because a product of commuting involutions has
finite order.

We first recover tetrahedra from this set.  Suppose a vector
subspace $U$ is contained in $\mathcal T$.  Any two coordinates
that occur somewhere in $U$ occur together in one vector of
$U$: choose a witness for each coordinate; if neither witness
has both coordinates, their sum does.  That vector is
supported on a clique.  It follows that the union of all
supports in $U$ is a clique.  Conversely, the full span of
every clique lies in $\mathcal T$.  Hence a maximal subspace
contained in $\mathcal T$ is exactly the full span of a
maximal clique.  These spaces are
\[
 V_\tau=\operatorname{span}_{\Ftwo}
       \{\mathbf e_u:u\text{ is a vertex of }\tau\},
\]
where $\tau$ is a maximal clique.  The triangulation is flag
and is a pure three-manifold triangulation, so its maximal
cliques are precisely its tetrahedra.

We next recover each coordinate line from those subspaces.
For a vertex $u$,
\begin{equation}\label{nor:line-intersection}
 \bigcap_{\tau\ni u}V_\tau=\langle\mathbf e_u\rangle,
\end{equation}
where the intersection ranges over tetrahedra containing $u$.
Indeed, if another vertex $v$ lay in all those tetrahedra,
then $v$ would occur in every maximal triangle of
$\lk_{\widehat{\mathcal L}}(u)$.  This link is a pure PL
two-sphere.  Every one of its simplices lies in a maximal
triangle, so the link would equal the star of $v$ and would be
a cone.  That contradicts its being a two-sphere.  This proves
\eqref{nor:line-intersection}.  Conversely, every
one-dimensional intersection of tetrahedron subspaces is a
coordinate line: an intersection of subspaces spanned by
subsets of a fixed basis is the span of the intersection of
those subsets.

An automorphism of $W$ preserves $\mathcal T$ and hence
permutes its maximal vector subspaces.  Its induced invertible
linear map preserves intersections, so
\eqref{nor:line-intersection} and the converse show that it
permutes all coordinate lines.  Over $\Ftwo$ it therefore
permutes the coordinate vectors themselves.  For distinct
$u,v$, adjacency is equivalent to
$\mathbf e_u+\mathbf e_v\in\mathcal T$.  The permutation
preserves adjacency, and flagness makes it a simplicial
automorphism.

The image of a Coxeter generator has finite order.  By
Lemma~\ref{nor:coxeter-facts} it is conjugate to a clique
product; its abelianization is the sum of the coordinates in
that product.  Since the abelianization just obtained is one
coordinate vector, the product consists of the corresponding
single generator.  Composing with the inverse generator
permutation gives the final assertion.
\end{proof}

It remains to make the conjugating elements for different
generators agree.  For right-angled Coxeter groups with finite
defining graphs,
comparison of conjugators on overlapping special subgroups
also appears in \cite[Section~4]{GutierrezKaul2008}.  Here we
compare stars and use connected links; this gives a proof for
the infinite nerve as well.

\begin{lemma}\label{nor:conjugating}
Let $X$ be $\widehat{\mathcal L}$ or the link in
$\widehat{\mathcal L}$ of a simplex of dimension $0$, $1$, or
$2$.  If an automorphism $\varphi$ of $W(X)$ sends each vertex
generator to a conjugate of itself, then $\varphi$ is inner.
\end{lemma}

\begin{proof}
The links in the statement are flag PL spheres of dimensions
$2$, $1$, and $0$.  We prove the assertion for them in
increasing dimension, and then for $\widehat{\mathcal L}$.
The link of a vertex at each induction step is one of the
lower-dimensional links already treated.

For the zero-sphere,
$W(X)=\langle r,s\mid r^2=s^2=1\rangle$ is infinite dihedral.
After an inner adjustment, arrange that $\varphi(r)=r$.
With $t=rs$, every conjugate of $s$, and hence its image,
has the form $t^j r$.
The subgroup generated by $r$ and $t^j r$ is
$\langle r,t^j\rangle$; it is the whole group only if
$j=\pm1$.  Since $s=t^{-1}r$ and $rsr=tr$, these two
possibilities give the identity or conjugation by $r$.
This proves the base case.

Now let $X$ have positive dimension and suppose the result is
known for its vertex links.  For each vertex $u$, we claim that
one element $g_u\in W(X)$ implements $\varphi$ on all the
generators of its star.  Choose $b_u$ with
$\varphi(u)=b_u u b_u^{-1}$ and put
$\psi=c_{b_u^{-1}}\circ\varphi$, where $c_g$ denotes
conjugation by $g$.  The automorphism $\psi$ fixes $u$ and
therefore preserves its centralizer.  Flagness identifies that
centralizer, by Lemma~\ref{nor:coxeter-facts}, as
\begin{equation}\label{nor:star-product}
 C_{W(X)}(u)
   =\langle u\rangle\times W(\lk_X(u)).
\end{equation}
For a link generator $v$, write $\psi(v)=q_v v q_v^{-1}$.
The image belongs to the star subgroup because it commutes
with $u$.  Apply the star retraction from
Lemma~\ref{ext:support} to this equation.  Writing the
retracted conjugator as $u^{\epsilon_v}k_v$ under
\eqref{nor:star-product}, with $k_v\in W(\lk_X(u))$, gives
\[
 \psi(v)=k_v v k_v^{-1}.
\]
Thus $\psi$ maps the link subgroup into itself and conjugates
each of its generators inside that subgroup.

This restricted map is onto.  Given $y\in W(\lk_X(u))$,
surjectivity of $\psi$ on the centralizer supplies a preimage
$u^\epsilon k$ with $k\in W(\lk_X(u))$.  Because $\psi(u)=u$
and $\psi(k)$ is in the link subgroup,
$y=u^\epsilon\psi(k)$.  Uniqueness of the direct-product
coordinates forces $\epsilon=0$.  Hence the restriction of
$\psi$ is a conjugating automorphism of the link subgroup.
By induction it is conjugation by some $k\in W(\lk_X(u))$.
That conjugation also fixes $u$, so $\psi$ is conjugation by
$k$ on the whole star.  Undoing the adjustment gives the
claimed element $g_u=b_uk$.

We now synchronize the elements $g_u$.  If $u,v$ are
adjacent, put $d_{uv}=g_u^{-1}g_v$.  The two star
conjugations agree on every generator in the intersection of
the stars, so $d_{uv}$ centralizes all those generators.
In particular it centralizes $u$ and $v$.  The centralizer
formula and the intersection formula of
Lemma~\ref{ext:support} first put it in the special subgroup
on the intersection of the star vertex sets.  Flagness
identifies the common neighbors with the vertices of the
edge link, giving
\[
 d_{uv}\in W_{\st_X(u)\cap\st_X(v)}
   =\langle u,v\rangle\times W(\lk_X(\{u,v\})).
\]
It is central in this subgroup.  The link factor has trivial
center.  If the edge link is empty this is immediate.
Otherwise it is a flag PL sphere.  A universal vertex in its
graph would extend every simplex to that vertex by flagness,
making the sphere a cone.  There is no such vertex, and the
center formula of Lemma~\ref{nor:coxeter-facts} applies.
Consequently
\begin{equation}\label{nor:edge-difference}
 d_{uv}\in\langle u,v\rangle,\qquad
 d_{uv}=u^{\lambda_{u,v}}v^{\lambda_{v,u}},
 \quad \lambda_{u,v},\lambda_{v,u}\in\Ftwo.
\end{equation}
The subgroup $\langle u,v\rangle$ is elementary abelian, so
$d_{vu}=d_{uv}^{-1}=d_{uv}$; the two coefficients in
\eqref{nor:edge-difference} depend on the underlying edge
and its indicated endpoint, not on its orientation.

For a triangle with vertices $u,v,z$, the exact identity
$d_{uv}d_{vz}d_{zu}=1$ follows by cancellation of the
$g$'s.  In the abelianization of $W(X)$, its $u$-coordinate
says
\[
 \lambda_{u,v}=\lambda_{u,z}.
\]
When $\dim X\geq2$, the link of every vertex is connected.
Any two neighbors of $u$ can therefore be joined by a link
edge path; each edge of that path gives such a triangle with
$u$.  The coefficient $\lambda_{u,v}$ is thus independent
of the neighbor $v$.

The remaining positive-dimensional case is a link circle.
Enumerate its vertices cyclically as
$u_0,\ldots,u_{m-1}$.  Cancellation of the $g$'s gives
\[
 d_{u_0u_1}d_{u_1u_2}\cdots d_{u_{m-1}u_0}=1.
\]
In its abelianization, the coordinate of $u_i$ occurs only
in the two incident edge differences.  Their coefficients
are equal.  This gives the same independence at each vertex
of the circle.

Write $\lambda_u$ for the common coefficient at $u$, and
replace $g_u$ by $g'_u=g_u u^{\lambda_u}$.  This still
implements the star conjugation because $u$ centralizes its
star.  For every edge $uv$, Equation
\eqref{nor:edge-difference} gives
\[
 (g'_u)^{-1}g'_v
   =u^{\lambda_u}d_{uv}v^{\lambda_v}=1.
\]
The graph of $X$ is connected, so all $g'_u$ are equal.
This last step requires only a finite edge path between two
vertices and therefore also holds for the infinite graph
$\widehat{\mathcal L}$.  The common element conjugates every
generator as $\varphi$ does, proving that $\varphi$ is inner.
\end{proof}

\begin{proposition}\label{nor:rigidity}
Every automorphism of $W$ is a composite of an inner
automorphism and an automorphism induced by a simplicial
automorphism of $\widehat{\mathcal L}$.
\end{proposition}

\begin{proof}
Lemma~\ref{nor:generator-lines} gives the simplicial
permutation and reduces the remaining automorphism to one
that sends each generator to a conjugate of itself.
Lemma~\ref{nor:conjugating}, applied to
$X=\widehat{\mathcal L}$, makes that remaining automorphism
inner.  Moving an inner automorphism past a generator
permutation, if necessary, gives the stated order of the
two factors.
\end{proof}

\subsection{The normalizer calculation}

We now apply Proposition~\ref{nor:rigidity} in the normalizer
calculation, using the full-preimage description of $\mathcal E_2$.

\begin{proposition}\label{nor:normalizer}
Inside $\Out(\Pi)$,
\begin{equation}\label{nor:normalizer-equality}
 N_{\Out(\Pi)}(\mathsf A)
   =\mathsf A\times\langle\delta\rangle.
\end{equation}
\end{proposition}

\begin{proof}
Let $\omega\in N_{\Out(\Pi)}(\mathsf A)$ and choose a
representative $F\in\Aut(\Pi)$.  View
$\mathcal E_2\subseteq\Aut(\Pi)$ using
Proposition~\ref{ext:embedding}.  It is the full preimage
of $\mathsf A$, so conjugation by $F$ restricts to an
automorphism
\[
 \Phi\colon\mathcal E_2\longrightarrow\mathcal E_2.
\]
It preserves the inner copy of $\Pi$.  More precisely,
$\Phi(c_\gamma)=c_{F(\gamma)}$ for $\gamma\in\Pi$.
Since $\Pi$ is centerless, under its identification with
$\Inn(\Pi)$ this says that $\Phi|_\Pi=F$.

The subgroup $W$ is characteristic in $\mathcal E_2$.
Indeed, every finite-order element projects trivially to the
torsion-free group $H_2$, and so belongs to $W$.  Conversely,
$W$ is generated by its vertex involutions.  It is therefore
exactly the subgroup generated by all finite-order elements
of $\mathcal E_2$, a characteristic description.
Proposition~\ref{nor:rigidity} applies to $\Phi|_W$.
After composing $\Phi$ with conjugation by an element of
$W$, we may assume that this restriction is a simplicial
permutation $\sigma$ of $\widehat{\mathcal L}$.  This
adjustment preserves $\Pi$ because $\Pi$ is normal in
$\mathcal E$.

The adjusted automorphism induces an automorphism $\beta$
of $\mathcal E_2/\Pi=\mathsf A$.  For a vertex $u$ above
$s\in\mathcal V$, it satisfies
\begin{equation}\label{nor:fiber-labels}
 \beta(a_s)=a_{\overline{\sigma(u)}}.
\end{equation}
The right side is therefore the same for all vertices over
$s$.  Because the labels $a_s$ are distinct, $\sigma$ takes
each fiber of the vertex projection into one fiber.
Moreover, the target fibers for distinct $s$ are distinct
by injectivity of $\beta$.  The finite set $\mathcal V$
thus acquires a permutation $f$ with
$\overline{\sigma(u)}=f(\overline u)$; bijectivity of
$\sigma$ makes the maps of fibers onto.  Every simplex of
$\mathcal L$ has a lift, whose image under $\sigma$ projects
to the simplex on the $f$-images of its vertices.  Applying
the same argument to $\sigma^{-1}$ shows that $f$ is a
simplicial automorphism.

Let $\pi\colon\mathsf A_0\to\mathsf A$ be the label quotient,
and let $f_*$ permute the coordinate basis of $\mathsf A_0$.
Equation \eqref{nor:fiber-labels} on each $e_s$ says
$\beta\pi=\pi f_*$.  Taking kernels, and using the same
identity for the inverses, gives
$f_*(\mathsf C)=\mathsf C$.  Lemma~\ref{mod:rigid-relations} excludes
every nonidentity simplicial automorphism with this property.  Hence $f=1$.

It follows that $\sigma$ is a simplicial automorphism over
the identity of $\mathcal L$.  To see explicitly that it is
a deck transformation, choose one lifted vertex $u_0$.
Regularity of $\widehat S\to S$ supplies $h_0\in H$ with
$h_0u_0=\sigma(u_0)$.  Both $h_0^{-1}\sigma$ and
$\mathrm{id}_{\widehat S}$ are lifts of the covering
projection $\widehat S\to S$, and they agree at $u_0$.
Uniqueness of lifts on the connected domain makes them equal.
Thus $\sigma$ is the deck permutation given by $h_0$.

Both $\mathcal E_2$ and $\Pi$ are normal in $\mathcal E$.
We may consequently compose further with conjugation by
$(1,h_0^{-1})\in\mathcal E$ and assume that $\Phi$ fixes
$W$ pointwise.  For $h\in H_2$, write
$\Phi(1,h)=(d,h')$ with $d\in W$ and $h'\in H_2$.
Applying $\Phi$ to the conjugation action of $h$ on $W$
shows
\[
 c_d\circ h'=h\quad\text{as automorphisms of }W.
\]
Inner automorphisms act trivially on $W_{\mathrm{ab}}$,
so $h'$ and $h$ induce the same permutation of its
coordinate basis.  The deck action on vertices is
faithful, giving $h'=h$.  The displayed identity then
places $d$ in $Z(W)$.  This center is trivial by
Lemma~\ref{nor:coxeter-facts}: the vertex set is infinite
because the infinite group $H$ acts freely on it, while
local finiteness makes every vertex star finite, so there
is no universal vertex.  Hence $d=1$.  The adjusted
$\Phi$ fixes $H_2$ as well as $W$ and is the identity on
$\mathcal E_2$.

Undoing the adjustments, the original $\Phi$ is
conjugation on $\mathcal E_2$ by an element $e\in\mathcal E$.
For $\gamma\in\Pi$, its action on the inner copy of $\Pi$
then gives
$c_{F(\gamma)}=c_{e\gamma e^{-1}}$.  Centerlessness of
$\Pi$ implies $F(\gamma)=e\gamma e^{-1}$ for every
$\gamma$.  Thus $\omega$ is the class of $e$ and belongs
to $\mathcal E/\Pi=\mathsf A\times\langle\delta\rangle$.
The reverse inclusion in \eqref{nor:normalizer-equality}
holds because that direct product is abelian and contains
$\mathsf A$.
\end{proof}

\begin{proof}[Proof of Theorem~\ref{thm:pair}]
The specialization of Proposition~\ref{mod:models} to the
chosen labels gives closed connected topological
four-manifolds $M,N$ with homotopy equivalences to $Q_2$.
Since $Q_2$ is aspherical, both manifolds are aspherical.
Their common fundamental group $\Pi$ is word-hyperbolic by
Proposition~\ref{hyp:word-hyperbolic}.
It remains to exclude a homeomorphism between them.

Use the realization subgroups $R_M,R_N\leq\Out(\Pi)$
from Proposition~\ref{sym:criterion}, with the markings
of these two models.  The ambient group $\Out(\Pi)$ is
finite by Proposition~\ref{hyp:finite-out}.  The graph
groups $\mathsf A$ and $\mathsf B_\ell$ are $2$-subgroups
of the same order.  Proposition~\ref{ext:outer-actions}
gives
\[
 \mathsf A\leq R_M,\qquad \mathsf B_\ell\leq R_N,
\]
and Proposition~\ref{ext:delta} gives $\delta\notin R_M$.
The element $\delta$ is an involution and
$\mathsf A\times\langle\delta\rangle$ is an internal direct
product by Proposition~\ref{ext:embedding}; this is the
normalizer of $\mathsf A$ by
Proposition~\ref{nor:normalizer}.  Finally,
Proposition~\ref{ext:nonconjugate} gives the required
nonconjugacy.  All hypotheses of Proposition~\ref{sym:criterion}
hold with $X=M$, $Y=N$, and $\mathsf B=\mathsf B_\ell$.
That proposition proves that $M$ and $N$ are not homeomorphic.
\end{proof}
